% =========================================================================
% GUIA E REGRAS DE HIFENIZAÇÃO MANUAL (EXCEÇÕES E TERMOS ESTRANGEIROS)
% =========================================================================
% O LaTeX moderno (com os pacotes babel [brazilian] e fontenc [T1]) separa
% silabicamente e de forma automática praticamente 100% das palavras da língua 
% portuguesa com acentuação correta.
%
% Portanto, você NÃO precisa adicionar palavras comuns em português aqui.
%
% Este arquivo deve ser utilizado EXCLUSIVAMENTE para casos de exceção:
% 1. Termos técnicos em língua estrangeira (inglês) usados no texto;
% 2. Nomes próprios de softwares, modelos matemáticos ou marcas;
% 3. Alguma palavra muito específica da sua pesquisa que porventura caia
%    exatamente na margem direita e gere um aviso de "Overfull \hbox".
%
% COMO DECLARAR:
% Utilize \hyphenation{si-la-ba-1-si-la-ba-2-...} com todas as letras em 
% minúsculas e hifens comuns (-) indicando onde a quebra é permitida.
%
% DICA ALTERNATIVA:
% Para hifenizar uma palavra uma única vez diretamente no corpo do texto,
% basta inserir um hífen discricionário: ex.: pa\-la\-vra.
% =========================================================================

% Exemplos de termos técnicos e estrangeirismos comuns em pesquisas:
\hyphenation{bench-mark}
\hyphenation{bench-mark-ing}
\hyphenation{feed-back}
\hyphenation{hard-ware}
\hyphenation{soft-ware}
\hyphenation{trade-off}
\hyphenation{deep-learn-ing}

% Exemplos de softwares e ferramentas de modelagem/simulação de transportes:
\hyphenation{Postgre-SQL}
\hyphenation{Trans-CAD}
\hyphenation{Open-Street-Map}
